\documentclass[11pt,a4paper]{article}
\usepackage[utf8]{inputenc}
\usepackage[T1]{fontenc}
\usepackage{amsmath,amssymb,amsfonts}
\usepackage{physics}
\usepackage{geometry}
\usepackage{booktabs}
\usepackage{hyperref}
\usepackage{mathtools}

\geometry{margin=2.5cm}
\hypersetup{colorlinks=true,linkcolor=blue,citecolor=blue,urlcolor=blue}

\title{\textbf{Quantum Gravitational Dynamics}\\[4pt]
\Large Cosmological Dynamics of the Graviton Field}
\author{Romeo Matshaba}
\date{University of South Africa}

\begin{document}
\maketitle

\begin{abstract}
We derive the exact cosmological field equations for the graviton field \(\sigma_\mu\) from the full QGD action. In a homogeneous and isotropic universe, the system reduces to a single fourth‑order nonlinear ordinary differential equation for the temporal component \(\varphi(t)=\sigma_t(t)\). Linearising around a de Sitter background, we obtain four mode solutions: a constant, a decaying mode \(\propto e^{-3Ht}\), a rapidly growing ghost mode \(\propto e^{t/\ell_Q}\), and a rapidly decaying mode \(\propto e^{-t/\ell_Q}\). The nonlinear self‑interaction \(Q_t = 2\varphi\dot\varphi^2\) is shown to saturate the ghost growth when the amplitude reaches the Planck scale. Computing the energy‑momentum tensor correctly from the action yields an equation of state \(w=-1\) for the kinetic term, confirming that the graviton field behaves as a cosmological constant. However, the constant‑velocity solution \(\dot\varphi = v\) does not satisfy the full field equation; hence no exact de Sitter attractor exists. The quantum stiffness term \(\ell_Q^2\Box_g^2\sigma_\mu\) contributes to the stress tensor but its natural scale is \(M_{\mathrm{Pl}}^4\), far above the observed dark energy density. The origin of dark energy remains an open problem, possibly requiring a potential or a detailed understanding of the saturated ghost dynamics.
\end{abstract}

\tableofcontents

%=============================================================
\section{Introduction}
\label{sec:intro}
%=============================================================

The QGD framework posits that gravity emerges from the phase of a Dirac wavefunction, with the graviton field \(\sigma_\mu = \partial_\mu S/(mc)\) as the fundamental variable. The metric is a composite field constructed from \(\sigma_\mu\) via the master metric. In previous chapters, we derived the field equations for \(\sigma_\mu\) from the action and studied their linearised structure, including the appearance of a fourth‑order quantum stiffness term \(\ell_Q^2\Box_g^2\sigma_\mu\). This term introduces additional degrees of freedom and modifies the dynamics at short distances.

In this paper, we focus on the cosmological implications of the theory. Under the assumptions of homogeneity and isotropy, the graviton field reduces to a single time‑dependent component \(\varphi(t) = \sigma_t(t)\). We derive the exact equation of motion from the full action, including the self‑interaction \(Q_\mu\) and matter coupling. The resulting fourth‑order nonlinear ODE is then analysed. We find the linearised mode decomposition, identify a rapidly growing ghost mode, and show how nonlinearities can saturate its growth. We also compute the energy‑momentum tensor correctly and find that the kinetic term gives an equation of state \(w=-1\), exactly as in the original cosmology chapter. The constant‑velocity solution is shown not to satisfy the full equation, so the origin of dark energy must be sought elsewhere, perhaps in the quantum stiffness term or a potential.

%=============================================================
\section{The Action and Field Equations}
\label{sec:action}
%=============================================================

We start from the complete QGD action (Chapter 2, Eq. (10)):
\begin{equation}
S[\sigma] = \int d^4x \sqrt{-g(\sigma)} \left( \frac{R[g(\sigma)]}{16\pi G} + \frac{\hbar^2}{2M} g^{\mu\nu} \nabla_\mu \sigma^\alpha \nabla_\nu \sigma_\alpha + \mathcal{L}_{\mathrm{matter}}(\psi, g(\sigma)) \right),
\label{eq:action}
\end{equation}
where \(g_{\mu\nu}\) is the composite metric built from \(\sigma\) via the master formula. Varying with respect to \(\sigma_\alpha\) and using the chain rule, we obtain (see Chapter 3 for details)
\begin{equation}
\frac{\hbar^2}{M} \square_g \sigma^\alpha = \frac{1}{16\pi G} \left( G^{\mu\nu} - 8\pi G T_{\mathrm{tot}}^{\mu\nu} \right) \frac{\partial g_{\mu\nu}}{\partial \sigma_\alpha},
\label{eq:field}
\end{equation}
where \(T_{\mathrm{tot}}^{\mu\nu} = T_{\mathrm{matter}}^{\mu\nu} + T_\sigma^{\mu\nu}\) includes the stress tensor of the \(\sigma\)-field itself. In the master metric, the derivative \(\partial g_{\mu\nu}/\partial \sigma_\alpha\) contains terms linear in \(\sigma\) and first derivatives, leading to a complicated expression. However, for a homogeneous and isotropic universe, the metric reduces to the FRW form, and the field equation simplifies considerably.

%=============================================================
\section{Homogeneous Cosmology}
\label{sec:cosmo}
%=============================================================

We assume a flat FRW metric
\begin{equation}
ds^2 = dt^2 - a(t)^2 \, d\mathbf{x}^2,
\end{equation}
and a homogeneous graviton field
\begin{equation}
\sigma_\mu(t) = (\varphi(t), 0,0,0).
\end{equation}
The covariant derivatives acting on a scalar are
\begin{equation}
\square_g \varphi = \ddot\varphi + 3H\dot\varphi, \qquad H = \frac{\dot a}{a}.
\end{equation}

The master metric (Chapter 2, Eq. (5)) for this configuration gives \(g_{00} = 1 - \varphi^2\), \(g_{ij} = -a^2\delta_{ij}\), up to possible corrections from the quantum stiffness term in the metric itself. For the purpose of deriving the field equation, we use the full expression from Chapter 3, which includes the \(Q_\mu\), \(G_\mu\), and stiffness terms. For a homogeneous field, the geometric term \(G_\mu\) vanishes by symmetry, and the self‑interaction \(Q_\mu\) (Eq. (3.5)) reduces to
\begin{equation}
Q_t = 2\varphi \dot\varphi^2.
\end{equation}
The matter source \(T_\mu\) (Eq. (3.7)) becomes
\begin{equation}
T_t = \frac{1}{2} T^{\mu\nu} \sigma_\nu = \frac{1}{2} \rho_m \varphi,
\end{equation}
where \(\rho_m\) is the matter energy density. The quantum stiffness term \(\kappa\ell_Q^2\square_g^2\sigma_\mu\) (with \(\ell_Q^2 = G\hbar/c^3\)) gives, for a scalar,
\begin{align}
\square_g^2 \varphi &= \frac{1}{\sqrt{-g}} \partial_\mu \left( \sqrt{-g} g^{\mu\nu} \partial_\nu (\square_g \varphi) \right) \\
&= \frac{d^2}{dt^2} (\ddot\varphi + 3H\dot\varphi) + 3H \frac{d}{dt} (\ddot\varphi + 3H\dot\varphi) + \text{terms involving } \dot H.
\end{align}
Expanding, we obtain
\begin{equation}
\square_g^2 \varphi = \ddddot\varphi + 3H \dddot\varphi + 3\dot H \ddot\varphi + (3H^2 + 3\dot H) \dot\varphi. \tag{1}
\end{equation}

Putting everything together, the \(\alpha = t\) component of the master field equation (Chapter 3, Eq. (3.2)) yields
\begin{equation}
\boxed{ \ddot\varphi + 3H\dot\varphi = 2\varphi \dot\varphi^2 + \frac{1}{2} \rho_m \varphi + \kappa\ell_Q^2 \left[ \ddddot\varphi + 3H \dddot\varphi + 3\dot H \ddot\varphi + (3H^2 + 3\dot H) \dot\varphi \right]. } \tag{2}
\end{equation}
This is the exact fourth‑order nonlinear equation governing the cosmological evolution of \(\varphi\). It is the central result of this section.

%=============================================================
\section{Linear Analysis and Mode Decomposition}
\label{sec:linear}
%=============================================================

To understand the dynamics, we first consider the linear regime where \(\varphi\) is small and we neglect the nonlinear term \(2\varphi\dot\varphi^2\) and the matter coupling. We also assume a de Sitter background with constant \(H\). Equation (2) then reduces to
\begin{equation}
\ddot\varphi + 3H\dot\varphi - \kappa\ell_Q^2 \left( \ddddot\varphi + 3H \dddot\varphi \right) = 0. \tag{3}
\end{equation}
Define \(\psi = \ddot\varphi + 3H\dot\varphi\). Then (3) becomes
\begin{equation}
\psi - \kappa\ell_Q^2 \left( \ddot\psi + 3H\dot\psi \right) = 0.
\end{equation}
This is a second‑order linear ODE for \(\psi\) with constant coefficients. Seeking solutions \(\psi \propto e^{rt}\), we obtain the characteristic equation
\begin{equation}
1 - \kappa\ell_Q^2 (r^2 + 3H r) = 0 \quad \Rightarrow \quad r^2 + 3H r - \frac{1}{\kappa\ell_Q^2} = 0.
\end{equation}
The roots are
\begin{equation}
r_\pm = \frac{-3H \pm \sqrt{9H^2 + 4/(\kappa\ell_Q^2)}}{2}.
\end{equation}
Since \(1/(\kappa\ell_Q^2) \sim M_{\mathrm{Pl}}^2\) is enormous compared to \(H^2\), we have approximately
\begin{equation}
r_+ \approx \frac{1}{\sqrt{\kappa}\,\ell_Q}, \qquad r_- \approx -\frac{1}{\sqrt{\kappa}\,\ell_Q} - \frac{3H}{2}.
\end{equation}
Thus
\begin{equation}
\psi(t) = A e^{r_+ t} + B e^{r_- t}.
\end{equation}
Now \(\varphi\) satisfies \(\ddot\varphi + 3H\dot\varphi = \psi\). The homogeneous solution of this equation is \(\varphi_h = C + D e^{-3Ht}\). Particular solutions for the exponentials are obtained by assuming \(\varphi \propto e^{r t}\), which gives
\begin{equation}
\varphi_p = \frac{\psi}{r^2 + 3H r}.
\end{equation}
Therefore the general solution for \(\varphi\) is
\begin{equation}
\varphi(t) = C + D e^{-3Ht} + \frac{A}{r_+^2 + 3H r_+} e^{r_+ t} + \frac{B}{r_-^2 + 3H r_-} e^{r_- t}. \tag{4}
\end{equation}
The four independent modes are:
\begin{itemize}
  \item a constant mode \(C\),
  \item a decaying mode \(e^{-3Ht}\) (timescale \(\sim 1/(3H)\)),
  \item a rapidly growing mode \(e^{r_+ t}\) with \(r_+ \approx 1/(\sqrt{\kappa}\,\ell_Q)\),
  \item a rapidly decaying mode \(e^{r_- t}\) with \(r_- \approx -1/(\sqrt{\kappa}\,\ell_Q)\).
\end{itemize}
The growing mode is the Ostrogradsky ghost inherent in fourth‑order theories. Its time scale is the Planck time \(\ell_Q/c \sim 10^{-43}\) s, so it would destroy the classical background almost instantaneously unless it is suppressed by initial conditions or saturated by nonlinearities.

%=============================================================
\section{Nonlinear Saturation of the Ghost}
\label{sec:nonlinear}
%=============================================================

The nonlinear term \(2\varphi\dot\varphi^2\) in Eq. (2) becomes important when the amplitude grows. To estimate the saturation, assume that the ghost mode dominates and write \(\varphi \sim \Phi(t) e^{r_+ t}\) with slowly varying \(\Phi\). Then \(\dot\varphi \sim r_+ \varphi\), \(\ddot\varphi \sim r_+^2 \varphi\), and the left‑hand side of (2) is dominated by the \(\kappa\ell_Q^2\) term:
\begin{equation}
\kappa\ell_Q^2 \ddddot\varphi \sim \kappa\ell_Q^2 r_+^4 \varphi \sim \frac{1}{\kappa\ell_Q^2} \varphi,
\end{equation}
since \(r_+^2 \approx 1/(\kappa\ell_Q^2)\). The right‑hand side contains the nonlinear term
\begin{equation}
2\varphi\dot\varphi^2 \sim 2\varphi (r_+^2 \varphi^2) = 2 r_+^2 \varphi^3 \sim \frac{2}{\kappa\ell_Q^2} \varphi^3.
\end{equation}
Thus the nonlinear term is of order \(\varphi^3\) while the linear term is of order \(\varphi\). Balancing them gives \(\varphi^2 \sim 1\), i.e., \(\varphi\) saturates at Planckian values \(\varphi \sim 1\) in Planck units. This suggests that the ghost growth is halted when the field amplitude reaches the Planck scale, at which point the nonlinearities become comparable to the linear terms. A full numerical integration would be required to confirm this picture and to determine the post‑saturation dynamics.

%=============================================================
\section{Energy-Momentum Tensor and Equation of State}
\label{sec:energy}
%=============================================================

The energy‑momentum tensor for the \(\sigma\)-field is obtained by varying the \(\sigma\)-kinetic term with respect to the metric. For the Lagrangian density \(\mathcal{L}_\sigma = \frac{\hbar^2}{2M} g^{\mu\nu} \nabla_\mu \sigma^\alpha \nabla_\nu \sigma_\alpha\), the stress tensor is given by
\begin{equation}
T_{\mu\nu}^{\sigma} = \frac{\hbar^2}{M} \left( \nabla_\mu \sigma^\alpha \nabla_\nu \sigma_\alpha - \frac{1}{2} g_{\mu\nu} \nabla_\lambda \sigma^\alpha \nabla^\lambda \sigma_\alpha \right).
\end{equation}
For a homogeneous field with only \(\varphi(t)\), we have \(\nabla_0 \sigma^0 = \dot\varphi\) and all other derivatives vanish. Then
\begin{align}
\nabla_0 \sigma^0 \nabla_0 \sigma_0 &= \dot\varphi^2, \\
\nabla_\lambda \sigma^\alpha \nabla^\lambda \sigma_\alpha &= g^{00} \nabla_0 \sigma^\alpha \nabla_0 \sigma_\alpha = \dot\varphi^2.
\end{align}
Thus
\begin{align}
T_{00}^{\sigma} &= \frac{\hbar^2}{M} \left( \dot\varphi^2 - \frac{1}{2} g_{00} \dot\varphi^2 \right) = \frac{\hbar^2}{2M} \dot\varphi^2, \\
T_{ij}^{\sigma} &= \frac{\hbar^2}{M} \left( 0 - \frac{1}{2} g_{ij} \dot\varphi^2 \right) = -\frac{\hbar^2}{2M} g_{ij} \dot\varphi^2.
\end{align}
Raising indices (with \(g^{00}=1\), \(g^{ij} = -\delta^{ij}/a^2\)), we obtain
\begin{align}
T^{00}_{\sigma} &= \frac{\hbar^2}{2M} \dot\varphi^2, \\
T^{ij}_{\sigma} &= \frac{\hbar^2}{2M a^2} \dot\varphi^2 \delta^{ij}.
\end{align}
For a perfect fluid, \(T^{\mu\nu} = (\rho + p) u^\mu u^\nu - p g^{\mu\nu}\). With \(u^\mu = (1,0,0,0)\), we have \(T^{00} = \rho\), \(T^{ij} = p g^{ij}\). Since \(g^{ij} = -\delta^{ij}/a^2\), this means \(T^{ij} = -p \delta^{ij}/a^2\). Equating with our expression gives
\begin{equation}
\frac{\hbar^2}{2M a^2} \dot\varphi^2 \delta^{ij} = -p \frac{\delta^{ij}}{a^2} \quad\Longrightarrow\quad p = -\frac{\hbar^2}{2M} \dot\varphi^2.
\end{equation}
Therefore
\begin{equation}
\rho_\sigma = \frac{\hbar^2}{2M} \dot\varphi^2, \qquad p_\sigma = -\frac{\hbar^2}{2M} \dot\varphi^2,
\end{equation}
and the equation of state is
\begin{equation}
\boxed{w_\sigma = \frac{p_\sigma}{\rho_\sigma} = -1},
\end{equation}
exactly, independent of the value of \(\dot\varphi\). This confirms the result of the original cosmology chapter: the kinetic term of the graviton field behaves as a cosmological constant.

%=============================================================
\section{Implications for Dark Energy}
\label{sec:DE}
%=============================================================

The constant‑velocity solution \(\dot\varphi = v\), if it existed, would give \(\rho_\sigma = \frac{\hbar^2}{2M} v^2 = \text{const}\) and \(p_\sigma = -\rho_\sigma\), exactly matching a cosmological constant. However, such a solution must satisfy the full field equation (2). Substituting \(\dot\varphi = v\) (so \(\ddot\varphi = 0\), etc.) into (2) with \(\rho_m=0\) and neglecting the stiffness term gives
\begin{equation}
3H v = 2\varphi v^2.
\end{equation}
Since \(\varphi\) grows linearly with time (\(\varphi = \varphi_0 + vt\)), this cannot hold for all \(t\) unless \(v=0\) or \(\varphi\) is constant (which would imply \(v=0\)). Thus no exact constant‑velocity solution exists in the full nonlinear theory.

Nevertheless, the Friedmann equation in the presence of the \(\sigma\)-field is
\begin{equation}
H^2 = \frac{8\pi G}{3c^2} \left( \rho_m + \frac{\hbar^2}{2M} \dot\varphi^2 \right).
\end{equation}
If \(\dot\varphi\) varies slowly, the \(\sigma\)-field can mimic a slowly varying dark energy component with \(w \approx -1\). The observed value \(\rho_\Lambda \approx 6\times10^{-10}\,\mathrm{J/m^3}\) would require \(\dot\varphi \sim \sqrt{2M\rho_\Lambda}/\hbar\). For \(M\) of order the Planck mass, this gives a tiny \(\dot\varphi\), consistent with slow roll.

The quantum stiffness term \(\ell_Q^2\Box_g^2\sigma_\mu\) contributes to the stress tensor at order \(\ell_Q^2\). Its natural scale is \(M_{\mathrm{Pl}}^4\), so it would dominate unless fine‑tuned. However, if the ghost mode saturates at Planckian amplitudes, the effective energy density could be of order \(M_{\mathrm{Pl}}^4\) during the quantum phase, but this would be unacceptable at late times. A more detailed analysis of the nonlinear saturation is required to determine whether the field can settle into a state with \(\dot\varphi\) small enough to yield the observed \(\rho_\Lambda\).

The complete coupled system (2) and the Friedmann equation must be solved numerically to determine the late‑time behavior. The presence of the ghost mode and its nonlinear saturation may lead to a natural mechanism for setting the scale of dark energy, but this remains an open question.

%=============================================================
\section{Conclusions}
\label{sec:conclusions}
%=============================================================

We have derived the exact cosmological field equation for the graviton field \(\varphi = \sigma_t\) from the QGD action. The equation is fourth‑order and nonlinear, incorporating self‑interaction, matter coupling, and quantum stiffness. Linearising around a de Sitter background, we identified four modes, including a rapidly growing ghost mode. We argued that nonlinearities may saturate this growth at Planckian amplitudes. The energy‑momentum tensor for the \(\sigma\)-field yields an equation of state \(w=-1\) for the kinetic term, confirming the cosmology chapter's result. However, the constant‑velocity solution does not satisfy the full equation, so no exact de Sitter attractor exists. The observed dark energy density could still arise if \(\dot\varphi\) is slowly varying, but a full numerical integration of (2) is required to determine whether such a slow‑roll solution exists and is an attractor. The quantum stiffness term may also play a role, but its natural scale is far too large, suggesting that a potential or a more subtle mechanism is needed. Future work should focus on numerical analysis and phase‑space methods to uncover the late‑time behavior of the system.

\end{document}
